\section{Introduction}
\label{sec:intro}

We consider the problem of simultaneously reconstructing \emph{4D scene geometry}
and estimating \emph{dense point motion} from a monocular RGB video of a dynamic scene 
in a feed-forward manner.
This formulation mirrors how the physical world operates: an object is structured by its geometry in 3D space,
as well as its motion across time.
Achieving this goal is highly challenging, as monocular 4D reconstruction is inherently ill-posed and dense temporal correspondences remain difficult,
especially under occlusions and significant motion.
However, a successful solution would have a wide spectrum of applications,
from video understanding to robotics~\cite{wang2017deepvo,qin2018vins} and world models~\cite{ha2018world, hafner2025mastering}.

Traditional methods tackle this problem by finding the pixel correspondences over time,
and then iteratively optimizing a 3D mesh to fit the RGB(D) observations~\cite{zollhofer2014real,newcombe2015dynamicfusion,innmann2016volumedeform}.
However, they often produce noisy results limited by the sensor,
and need per-scene optimization, which is less generalizable.
In the deep learning era, this problem is typically divided into two sub-tasks:
dynamic geometry reconstruction~\cite{zhang24monst3r,wang2025continuous,jiang2025geo4d}
and correspondence estimation~\cite{teed2020raft,karaev2024cotracker}, although they are inherently related, both relying on pixel correspondence in multi-view geometry~\cite{hartley2003multiple}.

Recent feed-forward methods such as St4RTrack~\cite{st4rtrack2025},
Dynamic Point Maps~\cite{sucar2025dynamic},
and Stereo4D~\cite{jin2024stereo4d},
have emerged as promising alternatives to address this problem,
by extending the \emph{static} 3D reconstruction networks,
like DUSt3R~\cite{wang2024dust3r} and MASt3R~\cite{leroy2024grounding},
to dynamic scenes via target-timepoint map prediction.
Even so, these methods process only \emph{pairwise} frames at once and rely on post-optimization to align the results,
reducing their ability to capture long-range motion coherence.

In this paper,
we present \method,
a framework that leverages video generators to jointly reconstruct 4D geometry and estimate dense motion
for a long monocular video sequence,
in a feed-forward manner,
\emph{without any post-optimization}.
We achieve this by proposing a \emph{world-centric} 4D representation
that denotes the dynamic scene using a sequence of point maps~\cite{wang2024dust3r,wang2025vggt,jiang2025geo4d}
and the corresponding scene flow~\cite{horn1981determining}, both defined in the world coordinate system.
This representation is intuitive and effective:
by eliminating the camera-induced motion components,
static background points ideally exhibit zero flow in the system,
making it easier to learn the motion patterns of dynamic objects.
By comparison,
prior works~~\cite{st4rtrack2025,sucar2025dynamic,jin2024stereo4d} only predict the target time point maps,
paired with the reference frame,
and do not explicitly model dense motion throughout the whole video.
We therefore argue that, to fully understand a dynamic 3D scene, it is crucial to jointly model both dense geometry and motion in a shared coordinate system throughout the \emph{entire} video sequence.


Another challenge of this task is the lack of large-scale in-the-wild datasets with dense geometry and motion.
Following recent trends in leveraging pre-trained generative models for 3D~\cite{ke2024repurposing,lu2025matrix3d,jiang2025geo4d}, we do not train our model from scratch but start from a pre-trained video generator~\cite{blattmann2023stable}.
This strategy significantly alleviates the data scarcity issue,
as the generator is trained on billions of visual data.
Moreover, the video generator inherently models spatiotemporal consistency across multiple frames,
making it well-suited for capturing long-term motion correspondence.
While Geo4D~\cite{jiang2025geo4d} has explored leveraging video generators~\cite{blattmann2023stable,xu2025geometrycrafter} for 4D reconstruction,
it only outputs \emph{independent} point maps for each frame,
without modeling dense motion across these points.

In this work, we take a further step towards jointly modeling dense geometry and motion.
We do so by encoding a unified 4D representation, combining point maps and scene flows, into a compact latent space.
Without the need to build cost volumes~\cite{teed2021raft} or establish dense correspondence~\cite{sucar2025dynamic} in pixel space, this integrated representation efficiently transfers the strong priors of the video generator to dense 4D geometry and motion reconstruction.


Moreover, we show that it is not necessary to strictly align the value range of 4D data to that of the original VAE in the diffusion model.
It is commonly believed that such alignment is essential for effectively leveraging pre-trained priors,
even for 3D geometry~\cite{ke2024repurposing,zhang2024world,lu2025matrix3d,jiang2025geo4d},
whose distribution differs substantially from that of natural images.
Our results, however, suggest otherwise.
Specifically, we adopt the canonical normalization for point maps,
\ie, centering the 3D coordinates and scaling them according to the scene's mean scale.
Despite this misalignment with the original RGB distribution in the VAE, \method{} still achieves strong generalization and accurate 4D reconstruction and motion estimation.
This finding challenges conventional beliefs and opens up new possibilities for leveraging diffusion models for geometric tasks.

To summarize, our key contributions are:
(1) We present \method,
a framework that leverages video generators to jointly reconstruct 4D scene geometry and estimate dense motion from monocular videos.
(2) We propose a novel 4D latent representation that unifies the modeling of geometry and motion, making our model simple but effective and easy to extend.
(3) We also show that strong generalization can be achieved without strictly aligning our 4D representation to the latent space of video diffusion,
challenging the conventional wisdom in diffusion-based 3D learning.


\begin{figure*}
    \centering
    \includegraphics[width=\linewidth]{figs/pipeline_figure.pdf}
    \caption{\textbf{Overview of \method}.
    % We define both scene geometry and motion within a shared world coordinate system. 
    We first train a novel \emph{4D VAE} (bottom-right), consisting of a \emph{Geometry VAE} and a \emph{Motion VAE}. 
    These two components jointly encode the point map and scene flow into a unified 4D latent representation. 
    Within the Diffusion Unet, we leverage the pretrained VAE from SVD (Stable Video Diffusion) to encode video latents as conditional inputs, 
    which are then channel-wise concatenated with our 4D latent to guide the denoising process. We only add noise to the 4D latents during model training for the Diffusion version.
    Note that we do not enforce the 4D latent distribution to strictly align with the original SVD VAE latent distribution.  
    And we find that this relaxed training strategy consistently improves the generalization performance of both the VAE and the Diffusion Unet.
    }
    \label{fig:overview}
    \vspace{-1em}
\end{figure*}